\section{A planar surface cannot have one nontrivial boundary}
\label{sec:diagram}

We now combine the two unitary lemmas. Corner-labelled relative diagrams
are classical; see Howie \cite[Section 1, Lemma 1]{Howie1983}.
Diagram obstructions and Euler-characteristic estimates also appear in
\cite[preprint Sections 3--4, especially Lemma 2]{KlyachkoLurye2012} and
\cite[preprint Sections 3--4]{Chen2026}. We give the disk-and-band argument in full,
using the spectral-phase comparison for its central estimate.
Fix a group $A$ and a cyclic word
\begin{equation}\label{eq:word}
 w=t^{s_1}a_1\cdots t^{s_n}a_n,
 \qquad a_i\in A,\quad s_i\in\{1,-1\},\quad
 \sum_{i=1}^n s_i=\pm1.
\end{equation}
Thus $n$ is odd and positive. The coefficients may equal $1$.
Indices on this word are taken modulo $n$.

A \emph{word disk} is an oriented disk whose boundary has $n$ disjoint
marked intervals, called slots, separated by labelled arcs, called corners.
A positive disk reads $w$ around its oriented boundary, and a negative disk
reads $w^{-1}$. Each slot is indexed by the occurrence of $t$ in
\eqref{eq:word} from which it comes. Its \emph{traversal sign} is $s_i$ on
a positive disk and $-s_i$ on a negative disk. The sign of a disk and the
traversal sign of a slot are different data.

Attach rectangular bands to pairs of slots, compatibly with orientation,
using every slot exactly once. Require each band to join opposite traversal
signs. The resulting compact oriented surface has boundary words in $A$:
read the corner labels in the boundary orientation, ignoring band sides.
Changing the starting point conjugates the word, which does not affect
whether it represents $1$.

\begin{theorem}[Planar boundary theorem]\label{thm:planar}
Let $\Sigma$ be a connected surface of genus zero constructed from a
nonempty finite collection of word disks for \eqref{eq:word} and bands as
above. If every boundary word except possibly one is trivial in $A$, then
the remaining boundary word is trivial as well.
\end{theorem}

\subsection{Slots, corners, and boundary cycles}

The total of all traversal signs is zero, since the bands pair opposite
signs. A positive disk contributes $\sum_i s_i$ and a negative disk its
negative. There are therefore $k$ disks of each kind, with $k\geq1$.
Let $D$ be the set of their $2nk$ slots. Define two permutations of $D$:
$\sigma$ takes each slot to the next slot on its oriented disk, and
$\theta$ pairs the two slots of a band. Write $c_d\in A$ for the corner
from $d$ to $\sigma d$. Thus
\begin{equation}\label{eq:corners}
 \begin{array}{c|c|c}
 \text{disk kind}&\text{successor}&\text{corner label}\\ \hline
 +&i\longmapsto i+1&a_i\\
 -&i+1\longmapsto i&a_i^{-1}.
 \end{array}
\end{equation}

The boundary circles are the cycles of $\sigma\theta$. Indeed, from the
start of a slot $d$, a boundary path follows a band side to the finish of
$\theta d$, then follows its corner to the start of $\sigma\theta d$.
For a cycle $z=(d_1,\ldots,d_{d_z})$ in this order, its word is
\begin{equation}\label{eq:boundaryword}
 r_z=c_{\theta d_1}\cdots c_{\theta d_{d_z}}.
\end{equation}
Figure~\ref{fig:boundary} records this local convention.

\begin{figure}[ht]
\centering
\begin{tikzpicture}[>=Stealth,scale=.94,every node/.style={font=\small}]
 \fill[gray!12] (0,0) rectangle (2.5,1.6);
 \draw (0,0)--(0,1.6) (2.5,0)--(2.5,1.6);
 \draw[very thick] (0,0)--(2.5,0) (0,1.6)--(2.5,1.6);
 \draw[->] (1.6,0)--(.8,0);
 \draw[->] (.8,1.6)--(1.6,1.6);
 \node[below] at (1.25,0) {$d$};
 \node[above] at (1.25,1.6) {$\theta d$};
 \node at (1.25,.8) {band};
 \node[below right] at (2.5,0) {start};
 \node[right] at (2.6,1.3) {finish};
 \draw[blue!65!black,very thick,->] (2.5,.2)--(2.5,1.4);
 \draw[blue!65!black,very thick,->] (2.5,1.6) to[out=45,in=180] (4.9,2.1);
 \node[above] at (4,2.15) {$c_{\theta d}$};
 \draw[very thick,->] (4.9,2.1)--(6.25,2.1);
 \node[below] at (5.65,2.1) {$\sigma\theta d$};
 \node[align=center] at (5,.55) {one boundary step:\\$d\ \xrightarrow{\theta}\ \theta d
 \ \xrightarrow{\sigma}\ \sigma\theta d$};
\end{tikzpicture}
\caption{The slot arrows follow the oriented boundaries of the word disks.
The blue path follows the boundary of the attached surface. The paired
slots have opposite traversal signs; the corner encountered after the band
is $c_{\theta d}$.}
\label{fig:boundary}
\end{figure}

There are $2k$ disks and $nk$ bands. If $q$ is the number of boundary
circles, the Euler characteristic gives
\begin{equation}\label{eq:euler}
 2-q=2k-nk,\qquad q=nk-2k+2.
\end{equation}
In particular $q\geq1$: the surface has nonempty boundary along its corners.

We will assign a small real phase $\eta_z$ to each boundary circle, with
$\sum_z\eta_z=0$. The phases will let us approach the spectral cut from
opposite sides on the trivial and possibly nontrivial boundary circles.
They can be realized by numbers $\gamma_d\in\R$
satisfying
\begin{equation}\label{eq:flows}
 \gamma_{\theta d}=-\gamma_d,\qquad
 \sum_{d\in z}\gamma_d=\eta_z.
\end{equation}
To see this, form a graph whose vertices are the cycles of $\sigma\theta$
and whose edges are the pairs $\{d,\theta d\}$. The graph is connected:
connectedness of $\Sigma$ makes $\langle\sigma,\theta\rangle$ transitive on
$D$, and this group equals $\langle\sigma\theta,\theta\rangle$.
On a spanning tree, any assignment of total sum zero is the divergence of
a real edge flow. For completeness, remove a leaf, put its required value
on its end of the incident edge and its negative on the other end, and
adjust the value required at the adjacent vertex. The final vertex has
value zero. Set all other edge flows to zero, including loops. This gives
\eqref{eq:flows}; if $q=1$, take every $\gamma_d=0$.

\subsection{An upper bound obtained by mixing the disks}

Use the algebra $M_{2nk}(\mathcal M)$ and its trace from
Section~\ref{sec:phase}, with matrix coordinates indexed by $D$.
On $\C^D\otimes\ell^2(A)$ define
\begin{align}\label{eq:shifts}
 S(e_d\otimes v)&=e_{\sigma d}\otimes L(c_d^{-1})v, &
 \Theta(e_d\otimes v)&=e^{2\pi i\gamma_d}e_{\theta d}\otimes v.
\end{align}
Both operators are unitary; \eqref{eq:flows} gives $\Theta^2=I$.
The product $S\Theta$ follows the boundary permutation $\sigma\theta$.
Its coefficient weights act on the left, so a complete circuit multiplies
the inverse corner labels in reverse order, giving the inverse boundary word.
We regard scalar matrices on $\C^D$ as elements of this algebra by tensoring
with the identity. For a scalar orthogonal projection $Q$,
$\tau_{2nk}(Q\otimes I)=\operatorname{rank}Q$, since $\tau(I)=1$.

Number the disks of each kind by $1,\ldots,k$. At each slot index $i$,
write $e_{+,i}(u)$ and $e_{-,i}(u)$ for the positive- and negative-disk
arrays with coefficient vector $u\in\C^k$. Given $X\in\U(k)$, put
\begin{equation}\label{eq:mixing}
 J e_{+,i}(u)=e_{-,i}(Xu),\qquad
 J e_{-,i}(u)=e_{+,i}(X^{-1}u).
\end{equation}
This scalar operator is a unitary involution. The corner convention gives
\begin{align*}
 S(e_{+,i}(u)\otimes v)&=e_{+,i+1}(u)\otimes L(a_i^{-1})v,\\
 S(e_{-,i}(u)\otimes v)&=e_{-,i-1}(u)\otimes L(a_{i-1})v.
\end{align*}
The same coefficient weight acts on all $k$ disks of a given kind.
Consequently $JSJ=S^{-1}$: on a positive array both sides shift to $i-1$
with weight $L(a_{i-1})$, and on a negative array both shift to $i+1$
with weight $L(a_i^{-1})$. Thus $SJ$ is a self-adjoint unitary. It exchanges
disk kinds, so its trace is zero. Its negative spectral projection
$(I-SJ)/2$ has trace $nk$, and hence
\begin{equation}\label{eq:phaseSJ}
 \ell(SJ)=\frac{nk}{2}.
\end{equation}

We now choose $X$ to make $J\Theta$ have a large fixed space. Let
$V_+$ and $V_-$ be the scalar coordinate spaces of slots with positive
and negative traversal signs. Each has dimension $nk$. Both $J$ and
$\Theta$ exchange them. In $V_+$, the block at index $i$ consists of
positive disks when $s_i=1$ and negative disks when $s_i=-1$; in $V_-$,
use the opposite kind. Order each block by disk number. In these bases
$\Theta:V_+\to V_-$ is a fixed matrix $P\in\U(nk)$, while $J:V_-\to V_+$
is $\diag(X^{-s_1},\ldots,X^{-s_n})$. Thus
\begin{equation}\label{eq:traversalblock}
 (J\Theta)|_{V_+}=\diag(X^{-s_1},\ldots,X^{-s_n})P.
\end{equation}
Lemma~\ref{lem:fixed}, applied to the signs $-s_i$, supplies $X$ for which
this block has at least $k$ fixed directions. Moreover,
$J(J\Theta)J=(J\Theta)^{-1}$, and $J$ exchanges $V_+$ and $V_-$.
The other block therefore also has at least $k$ fixed directions.

The eigenvalues of the scalar unitary $J\Theta$ are paired with their
inverses, with multiplicities. Since
$h(\lambda)+h(\lambda^{-1})=1$ for $\lambda\ne1$ and both terms vanish
at $1$, we obtain
\begin{equation}\label{eq:phaseJTheta}
 \ell(J\Theta)=\frac{2nk-\dim\ker(J\Theta-I)}{2}\leq nk-k.
\end{equation}
Apply Lemma~\ref{lem:phase} to $S\Theta=(SJ)(J\Theta)$:
\begin{equation}\label{eq:upper}
 \ell(S\Theta)\leq\frac32 nk-k.
\end{equation}
This bound holds for every choice of the boundary phases in
\eqref{eq:flows}. The choice of $X$ may depend on those phases.

\subsection{Computing the phase along boundary circles}

We compute the same trace directly from the boundary cycles. On a cycle
$z=(d_1,\ldots,d_{d_z})$, one circuit of $S\Theta$ starting at $d_1$ has
weight
\[
 e^{2\pi i\eta_z}
 L(c_{\theta d_{d_z}}^{-1})\cdots L(c_{\theta d_1}^{-1})
 =e^{2\pi i\eta_z}L(r_z^{-1}).
\]
Circuit weights at other starting points are unitary conjugates.
If $T_z$ is the restriction of $S\Theta$ to this cycle, then
\begin{equation}\label{eq:cyclephase}
 \ell(T_z)=\frac{d_z-1}{2}
   +\tau\bigl(h(e^{2\pi i\eta_z}L(r_z^{-1}))\bigr).
\end{equation}
Here and below $\ell$ uses the matrix trace of the appropriate size.
To prove the formula, let $d=d_z$ and $\zeta=e^{2\pi i/d}$.
Rephasing the successive coordinate vectors by $1,\zeta,\ldots,\zeta^{d-1}$
conjugates $T_z$ to $\zeta T_z$. For every $\lambda\in\T$ the exact identity
\[
 \sum_{j=0}^{d-1}h(\zeta^j\lambda)
   =\frac{d-1}{2}+h(\lambda^d)
\]
follows by listing the $d$ fractional parts; it also holds when
$\lambda^d=1$, because $h(1)=0$. Apply bounded Borel functional calculus
and take the trace. Each term on the left has trace $\ell(T_z)$.
The constant on the right has trace $d(d-1)/2$, and $T_z^d$ is diagonal
with $d$ conjugate circuit weights. Division by $d$ gives
\eqref{eq:cyclephase}.

\begin{proof}[Proof of Theorem~\ref{thm:planar}]
Let $z_0$ be the possibly nontrivial boundary circle.
We approach the phase cut from below on the trivial circles and from
above on $z_0$. Thus each trivial word contributes a limit of one, while
the exceptional word retains its nonnegative phase. For $\epsilon>0$ put
\[
 \eta_z=-\epsilon\quad(z\ne z_0),\qquad
 \eta_{z_0}=(q-1)\epsilon.
\]
These phases sum to zero. On each trivial circle the last term of
\eqref{eq:cyclephase} tends to $1$ as $\epsilon\downarrow0$.
On $z_0$ it tends to $\tau(h(L(r_{z_0}^{-1})))$: the fractional part
converges to its prescribed value from the nonnegative-shift side,
including at $1$, and the trace obeys bounded convergence.
For $q=1$ the exceptional phase is identically zero and there are no other
circles, so the same statement holds.

Since $\sum_zd_z=2nk$, summing \eqref{eq:cyclephase} and then using
\eqref{eq:euler} gives
\begin{align*}
 \lim_{\epsilon\downarrow0}\ell(S\Theta)
 &=\frac{2nk-q}{2}+q-1+\tau(h(L(r_{z_0}^{-1})))\\
 &=\frac32 nk-k+\tau(h(L(r_{z_0}^{-1}))).
\end{align*}
The uniform bound \eqref{eq:upper} forces the last term to vanish.
Faithfulness of $\tau$ and positivity of $h$ imply
$h(L(r_{z_0}^{-1}))=0$, hence $L(r_{z_0}^{-1})=I$ and $r_{z_0}=1$.
Only the uniform bound is used; no limiting choice of $X$ or $J$ is needed.
\end{proof}
